\chapter{Meeting and Progress Evidence}
The team met the supervisor weekly, on Tuesdays, throughout the reporting period. The
table below records the sequence of those meetings by theme rather than listing each
occurrence separately, since the weekly cadence was uninterrupted.

\begin{longtable}{p{0.20\textwidth} p{0.30\textwidth} p{0.42\textwidth}}
\caption{Supervisor meetings and their outcomes}\label{tab:meetings}\\
\toprule
\textbf{Period} & \textbf{Focus} & \textbf{Outcome and where it is recorded} \\
\midrule
\endfirsthead
\toprule
\textbf{Period} & \textbf{Focus} & \textbf{Outcome and where it is recorded} \\
\midrule
\endhead
Before 11 September 2026 &
Proposal preparation and defence rehearsal &
Proposal defended on 11 September 2026 \\
\addlinespace
Weeks following the defence &
Review of the panel's comments, item by item &
Six feedback items recorded with the action taken against each, in the action
register (Section~\ref{sec:delays}) \\
\addlinespace
Subsequent weeks &
Strategy for the work after the defence: scope, attack classes, and the order in
which to establish the central claim &
The decision to test the cross-zone feature claim ahead of any federated
infrastructure, so that a null result would be found early (Section~\ref{sec:iteration-strategy}); the
scope position recorded in Section~\ref{sec:scope}\\
\addlinespace
Weeks before the mid evaluation &
Mid-evaluation preparation: evidence to present and gaps to close &
Iteration~1 results consolidated; remaining work sequenced in Chapter~\ref{ch:plan} \\
\bottomrule
\end{longtable}

\textbf{Assessment schedule.} The proposal was defended on 11 September 2026. The
mid evaluation falls in the week of 13 October 2026. The final FYP-1 evaluation is
scheduled for December 2026.

% [[FILL]] If your department requires per-meeting records rather than the thematic
% summary above, expand this into one row per Tuesday with the date, attendees,
% decisions taken, and action items with owner and status.
